% Analisi dei requisiti

\section{Introduzione}

Scopo della Prova Finale di Reti Logiche \`e implementare un encoder per Working Zones come descritto dalla specifica fornita. Per la realizzazione del progetto \`e stata utilizzata la fpga xc7a200tfbg484-1 nell'editor di Vivado Webpack 2019; la documentazione utilizzata per la realizzazione del progetto consiste nelle guide base di Vivado e di VHDL fornite dai docenti.

\subsection{Scelte progettuali}

Secondo la specifica definita, per la realizzazione del progetto sono stati utilizzati due registri per mantenere l'indirizzo da codificare e l'indirizzo della Working Zone analizzata al momento. Inoltre, si \`e deciso di mantenere un registro contatore per le operazioni effettuate, usato anche per comandare gli indirizzi cui accedere; in questo modo la gestione dell'accesso alla memoria \`e stato delegato al data path invece che alla macchina a stati finiti.

In base alle richieste sono anche state effettuate le seguenti scelte:
\begin{itemize}	
	\item Poich\'e non viene dichiarato che le WZ si trovino sempre in un ordine definito il circuito analizzer\`a tutti gli indirizzi prima di giungere alla conclusione di non appartenenza a nessuna WZ dell'indirizzo.
	
	\label{eccezione}\item Per quanto riguarda i valori non validi, il circuito dovr\`a considerare gli indirizzi di input come interi unsigned, quindi non sar\`a possibile inserire valori al di sotto di quelli considerati validi; se viene inserito un valore da codificare superiore a 127 il comportamento non dovr\`a differire da un'esecuzione normale, ma in caso di indirizzo non presente nelle WZ il MSB verr\`a sovrascritto dal WZ\_BIT, mantenendo quindi un comportamento consistente anche in casi eccezionali.
\end{itemize}

\subsection{Ottimizzazione con pipeline}

Per sfruttare al meglio la capacit\`a di parallelizzazione dell'hardware si \`e optato per una soluzione con una pipeline, cosicch\'e i moduli di gestione indirizzi, input e analisi (illustrati in seguito) potessero lavorare in contemporanea, riducendo notevolmente il tempo di esecuzione. A causa di questa scelta \`e stato per\`o necessario utilizzare un registro contatore per gestire l'offset tra i moduli.

Un esempio di esecuzione delle tre componenti in parallelo \`e mostrato in Figura \ref{pipeline}: una volta iniziata l'esecuzione il circuito richiede alla memoria prima l'indirizzo da codificare, quindi l'indirizzo della prima Working Zone. Da questo punto inizia l'esecuzione parallela a pieno regime, per cui avvengono contemporaneamente:
\begin{itemize}
	\item Analisi dell'appartenenza a una WZ;
	\item Caricamento della WZ successiva (se esiste);
	\item Richiesta di quella dopo ancora (se esiste).
\end{itemize}
Se l'analisi di appartenenza \`e positiva il sistema viene congelato e si procede all'output, altrimenti questo avviene una volta terminata l'analisi dell'ultima Working Zone.

\begin{figure}[t]
	\centering
	{\includegraphics[width=\textwidth,keepaspectratio]
		{pipeline.eps}}
	\caption{Esempio di funzionamento della pipeline per caricamento e analisi}
	\label{pipeline} 
\end{figure}